\chapter{Learning Theory — Concentration Bound}
%%% generated by the blueprint pipeline from TCSlib.LearningTheory.JohnsonLindenstrauss.ConcentrationBound
\section{Overview}

This module proves the single-vector concentration bound for the
Johnson--Lindenstrauss random projection.  A random matrix $A$ with i.i.d.\
$N(0,1/k)$ entries is shown to preserve the squared norm of any fixed vector $x$
up to a multiplicative $(1\pm\varepsilon)$ factor, with failure probability at
most $2\exp(-k\varepsilon^2/8)$.  The proof reduces to a centered chi-squared
moment-generating-function (MGF) bound obtained via a classical Taylor
inequality, then packages the result into the abstract Bernstein-MGF framework
so that the same tail machinery applies to other sub-Gaussian random matrix
families.

%%%%%%%%%%%%%%%%%%%%%%%%%%%%%%%%%%%%%%%%%%%%%%%%%%%%%%%%%%%%%%%%%%%%%%%%%%%%%%%
\section{Declarations}

\begin{definition}[Bad distortion event for a single vector]
\lean{BadSingle}
\label{BadSingle}
\leanok
For a parameter $\varepsilon\in\mathbb{R}$, a matrix $A$, and a vector
$x\in\mathbb{R}^d$, the predicate $\texttt{BadSingle}\,\varepsilon\,A\,x$ holds when
the squared-norm distortion is strictly too large:
\[
  \varepsilon \|x\|^2 < \bigl|\,\|Ax\|^2 - \|x\|^2\,\bigr|.
\]
This is the ``bad event'' whose probability must be controlled by the JL
concentration argument.
\end{definition}

\begin{lemma}[Concentration bound at the zero vector]
\lean{concentration_zero}
\label{concentration_zero}
\leanok
\uses{BadSingle}
\difficulty{3}
Let $\mu$ be a probability measure on a measurable space $\Omega$, let $\omega \mapsto
A(\omega)$ be a family of $k \times d$ real matrices indexed by $\Omega$, and let
$\varepsilon \in \bbr$. Then the probability that the zero vector $0 \in \bbr^d$
triggers the bad distortion event, namely that
\[
  \varepsilon\,\|0\|^2 < \bigl|\,\|A(\omega)\,0\|^2 - \|0\|^2\,\bigr|,
\]
satisfies
\[
\mu\bigl\{\omega : \varepsilon\,\|0\|^2 < \bigl|\,\|A(\omega)\,0\|^2 -
\|0\|^2\,\bigr|\bigr\}
    \;\le\; 2\exp\!\Bigl(-\tfrac{k\,\varepsilon^2}{8}\Bigr).
\]
\end{lemma}

\begin{lemma}[Scaling a centered real Gaussian]
\lean{map_const_mul_gaussian}
\label{map_const_mul_gaussian}
\leanok
\difficulty{3}
Let $X\colon \Omega \to \bbr$ be a measurable random variable on a measure space
$(\Omega, \mu)$ whose law is the centered Gaussian $N(0, \sigma_0)$ with variance
$\sigma_0 \ge 0$. Then for every scalar $c \in \bbr$, the scaled variable $c\,X$ has law
\[
  N\!\bigl(0,\; c^2 \sigma_0\bigr),
\]
the centered Gaussian with variance $c^2 \sigma_0$.
\end{lemma}

\begin{lemma}[Weighted sum of independent Gaussians is Gaussian]
\lean{sum_scaled_iid_gaussian_map}
\label{sum_scaled_iid_gaussian_map}
\leanok
\uses{map_const_mul_gaussian}
\difficulty{5}
Let $(\Omega,\mu)$ be a probability space, and let $(Y_j)_{j}$ be a family of
independent, measurable real-valued random variables, each distributed as a centred
Gaussian with mean $0$ and variance $\sigma_0 \ge 0$. Then for any real coefficients
$(x_j)_j$ and any finite index set $s$, the random variable $\sum_{j \in s} x_j\, Y_j$
has distribution $N\!\Bigl(0,\; \bigl(\sum_{j \in s} x_j^2\bigr)\,\sigma_0\Bigr)$; that
is, its law is the centred Gaussian on $\bbr$ with variance $\sigma_0 \sum_{j \in s}
x_j^2$.
\end{lemma}

\begin{lemma}[Independence of coordinates of a random matrix–vector product]
\lean{rows_indep}
\label{rows_indep}
\leanok
\difficulty{3}
Let $(\Omega,\mu)$ be a probability space, and let $A$ be a random $k\times d$ real
matrix, that is, a measurable map assigning to each $\omega\in\Omega$ a matrix
$A(\omega)\in\bbr^{k\times d}$. Suppose the rows of $A$ are mutually independent,
meaning the $k$ random vectors
$\omega\mapsto\bigl(A(\omega)_{i,1},\dots,A(\omega)_{i,d}\bigr)\in\bbr^{d}$, for
$i=1,\dots,k$, are mutually independent under $\mu$. Then for every fixed $x\in\bbr^{d}$
the $k$ coordinates of the product $A(\omega)\,x$, namely the real-valued random
variables $\omega\mapsto (A(\omega)\,x)_i=\sum_{j=1}^{d}A(\omega)_{i,j}\,x_j$ for
$i=1,\dots,k$, are mutually independent under $\mu$.
\end{lemma}

\begin{lemma}[Coordinates of a matrix-vector product]
\lean{toEuclideanLin_apply_eq_sum}
\label{toEuclideanLin_apply_eq_sum}
\leanok
\difficulty{1}
Let $A\in\bbr^{k\times d}$ be a $k\times d$ real matrix, regarded as the linear map from
the Euclidean space $\bbr^{d}$ to $\bbr^{k}$ that it induces, and let $x\in\bbr^{d}$.
Then for each index $i\in\{1,\dots,k\}$, the $i$-th coordinate of the image $Ax$ is
given by
\[
  (Ax)_i \;=\; \sum_{j} A_{ij}\,x_j .
\]
\end{lemma}

\begin{lemma}[Squared Euclidean norm as a sum of squared coordinates]
\lean{norm_sq_toEuclideanLin}
\label{norm_sq_toEuclideanLin}
\leanok
\difficulty{2}
Let $A$ be a $k \times d$ real matrix, regarded as the linear map $x \mapsto Ax$ from
the Euclidean space $\bbr^d$ to $\bbr^k$, and let $x \in \bbr^d$. Then the squared
Euclidean norm of the image is the sum of the squares of its coordinates,
\[
  \norm{Ax}^2 = \sum_{i=1}^{k} (Ax)_i^2 .
\]
\end{lemma}

\begin{lemma}[Derivative of a Taylor auxiliary function]
\lean{taylorAux_hasDerivAt}
\label{taylorAux_hasDerivAt}
\leanok
\difficulty{4}
Let $h \colon \bbr \to \bbr$ be the function $h(x) = x^2 + x + \log(1 - x)$. Then at
every point $u < 1$, the function $h$ is differentiable with derivative
\[
h'(u) = \frac{u(1 - 2u)}{1 - u}.
\]
\end{lemma}

\begin{lemma}[An auxiliary expression vanishes at zero]
\lean{taylorAux_zero}
\label{taylorAux_zero}
\leanok
\difficulty{1}
The real-valued expression $x^2 + x + \log(1 - x)$ takes the value $0$ at $x = 0$; that
is, $0^2 + 0 + \log(1 - 0) = 0$.
\end{lemma}

\begin{lemma}[Taylor bound for the centered chi-squared log-MGF]
\lean{neg_log_one_sub_two_mul_le_two_sq}
\label{neg_log_one_sub_two_mul_le_two_sq}
\leanok
\uses{taylorAux_hasDerivAt, taylorAux_zero}
\difficulty{6}
For every real number $s$ with $\abs{s} \le 1/4$, one has
\[
  -s - \tfrac{1}{2}\log(1 - 2s) \;\le\; 2s^2 .
\]
\end{lemma}

\begin{lemma}[Quadratic Gaussian integral]
\lean{integral_exp_mul_sq_standardGaussian}
\label{integral_exp_mul_sq_standardGaussian}
\leanok
\difficulty{5}
Let $Z$ be a standard Gaussian random variable, $Z \sim N(0,1)$, and let $s \in \bbr$
satisfy $2s < 1$. Then
\[
  \E\!\left[e^{s Z^2}\right]
  = \int_{\bbr} e^{s z^2}\, \mathrm{d}N(0,1)(z)
  = \frac{1}{\sqrt{1 - 2s}}.
\]
\end{lemma}

\begin{lemma}[Quadratic exponential integral against a centered Gaussian]
\lean{integral_exp_mul_sq_gaussianReal_zero}
\label{integral_exp_mul_sq_gaussianReal_zero}
\leanok
\uses{integral_exp_mul_sq_standardGaussian}
\difficulty{5}
Let $v > 0$ and write $N(0,v)$ for the centered Gaussian distribution on $\bbr$ with
mean $0$ and variance $v$. Then for every real number $t$ with $2tv < 1$,
\[
  \int_{\bbr} e^{t y^2}\, \mathrm{d}N(0,v)(y) = \frac{1}{\sqrt{1 - 2tv}}.
\]
\end{lemma}

\begin{lemma}[Integrability of a Gaussian moment generating integrand]
\lean{integrable_exp_mul_sq_gaussianReal_zero}
\label{integrable_exp_mul_sq_gaussianReal_zero}
\leanok
\difficulty{4}
Let $v$ be a nonnegative real number with $v \ne 0$, and let $N(0,v)$ denote the centred
Gaussian probability measure on $\bbr$ with variance $v$. If $t \in \bbr$ satisfies $2tv
< 1$, then the function $y \mapsto e^{t y^2}$ is integrable with respect to $N(0,v)$.
\end{lemma}

\begin{theorem}[Centered chi-squared MGF bound]
\lean{centered_chi_squared_step}
\label{centered_chi_squared_step}
\leanok
\uses{integrable_exp_mul_sq_gaussianReal_zero, integral_exp_mul_sq_gaussianReal_zero, neg_log_one_sub_two_mul_le_two_sq}
\difficulty{6}
Let $(\Omega,\mu)$ be a probability space and let $k$ be a positive natural number.
Suppose $Y:\Omega\to\bbr$ is measurable and normally distributed with mean $0$ and
variance $1/k$, that is, its law $Y_*\mu$ equals $N(0,1/k)$. Then for every real $t$
with $\abs{t}\le k/4$, the function $\omega\mapsto \exp\!\bigl(t\,(Y(\omega)^2 -
1/k)\bigr)$ is $\mu$-integrable, and the moment generating function of the centered
variable $Y^2 - 1/k$ satisfies
\[
  \E_\mu\!\left[e^{t(Y^2 - 1/k)}\right] \;\le\; \exp\!\Bigl(\frac{2t^2}{k^2}\Bigr).
\]
\end{theorem}

\begin{theorem}[Bernstein MGF bound for a centered chi-squared summand]
\lean{hasBernsteinMGF_centered_chi_squared}
\label{hasBernsteinMGF_centered_chi_squared}
\leanok
\uses{ProbabilityTheory.HasBernsteinMGF, centered_chi_squared_step}
\difficulty{2}
Fix an integer $k \ge 1$, and let $Y$ be a real-valued random variable on a probability
space $(\Omega, \mu)$ whose law is the centered Gaussian $\mathcal{N}(0, 1/k)$ of mean
$0$ and variance $1/k$. Then the centered square $Y^2 - 1/k$ satisfies the Bernstein MGF
condition on $(\Omega, \mu)$ with quadratic-growth parameter $c = 2/k^2$ and radius
$t_{\max} = k/4$.
\end{theorem}

\begin{lemma}[Chi-squared tail bound for independent $N(0,1/k)$ variables]
\lean{chi_squared_tail}
\label{chi_squared_tail}
\leanok
\uses{ProbabilityTheory.HasBernsteinMGF, ProbabilityTheory.HasBernsteinMGF.measure_abs_gt_le, ProbabilityTheory.HasBernsteinMGF.sum_of_iIndepFun, hasBernsteinMGF_centered_chi_squared}
\difficulty{5}
Let $\mu$ be a probability measure on $\Omega$, let $k \ge 1$, and let $Y_1,\dots,Y_k :
\Omega \to \bbr$ be independent random variables, each with law $N(0, 1/k)$. Then for
every $\varepsilon$ with $0 < \varepsilon < 1$,
\[
  \mu\!\left\{\omega \;\middle|\; \varepsilon
    < \Bigl|\,\textstyle\sum_{i=1}^{k} Y_i(\omega)^2 - 1\Bigr|\right\}
  \;\le\; 2\exp\!\Bigl(-\tfrac{k\,\varepsilon^2}{8}\Bigr).
\]
\end{lemma}

\begin{theorem}[Johnson–Lindenstrauss concentration for a single vector]
\lean{jl_concentration_single_via_chi_squared}
\label{jl_concentration_single_via_chi_squared}
\leanok
\uses{BadSingle, chi_squared_tail, concentration_zero, map_const_mul_gaussian, norm_sq_toEuclideanLin, rows_indep, sum_scaled_iid_gaussian_map, toEuclideanLin_apply_eq_sum}
\difficulty{6}
Let $k \ge 1$, and let $A$ be a random $k \times d$ real matrix on a probability space
$(\Omega, \mu)$ whose $kd$ entries $A_{ij}$ are jointly independent, each Gaussian with
law $N(0, 1/k)$. Then for every fixed vector $x \in \bbr^d$ and every $\varepsilon$ with
$0 < \varepsilon < 1$,
\[
\mu\bigl\{\,\omega : \bigl|\,\|A(\omega)\,x\|^2 - \|x\|^2\,\bigr| > \varepsilon\,\|x\|^2
\,\bigr\}
  \;\le\; 2\exp\!\Bigl(-\frac{k\,\varepsilon^2}{8}\Bigr),
\]
where $\|A(\omega)\,x\|$ denotes the Euclidean norm of the image $A(\omega)\,x \in
\bbr^k$.
\end{theorem}

\begin{theorem}[Distribution-free Johnson–Lindenstrauss concentration via Bernstein tails]
\lean{jl_concentration_single_via_bernstein}
\label{jl_concentration_single_via_bernstein}
\leanok
\uses{ProbabilityTheory.HasBernsteinMGF, ProbabilityTheory.centered_squared_iid_tail, norm_sq_toEuclideanLin}
\difficulty{4}
Fix an integer $k \ge 1$, a dimension $d$, a probability space $(\Omega, \mu)$, and a
vector $x \in \bbr^d$. Let $A : \Omega \to \bbr^{k \times d}$ be a random matrix, and
for each outcome $\omega$ write $y(\omega) = A(\omega)\,x \in \bbr^k$ for the image of
$x$ under the random linear map, with coordinates $y_1(\omega), \dots, y_k(\omega)$.
Suppose each coordinate $y_i$ is measurable, the coordinates $y_1, \dots, y_k$ are
independent, and, for some constants $c > 0$ and $t_{\max}$, each centered squared
coordinate $y_i^2 - \norm{x}^2/k$ satisfies the Bernstein moment generating function
condition with parameters $(c, t_{\max})$. Then for every $s$ with $0 \le s \le
2kc\,t_{\max}$,
\[
  \mu\bigl\{\omega \;:\; s < \bigl|\,\norm{A(\omega)\,x}^2 - \norm{x}^2\,\bigr|\bigr\}
  \;\le\; 2\exp\!\Bigl(-\frac{s^2}{4kc}\Bigr).
\]
\end{theorem}
